% Derived analytic integration fragment, 2026-10-08.
% No numerical integration or new global Chen constant is certified here.
% Uses the existing prefix interface and the selected-subproduct repair.

\subsection{A safe ordinary upper option for four specified primes}
\label{sec:four-prefix-safe-minimum}

Fix $0<a<h<d<1/3$.  Write the four specified prime exponents in
decreasing order as $t=(x,y,z,w)$ and put
\[
 S=x+y+z+w,\quad
 W(t)=\max\{2x,x+2y,x+y+2z,x+y+z+2w\},
 \qquad d\mu(t)=\frac{dx\,dy\,dz\,dw}{xyzw}.
\]
The two tuple domains, with the full fourth-prime range, are
\[
 \Omega_{10}=\{a<w<z<y<x<h\},\qquad
 \Omega_{11}=\{a<w<z<y<h<x<d\}.
\]
The quotient in each original count is sifted at the second smallest
specified prime, whose exponent is $z$.  The four specified primes
need not be the four largest prime factors of that quotient count.
They will, however, be the four largest primes in every modulus atom
of the ordinary upper sieve constructed below.  These are different
statements, and only the latter is used in factorization.

Put
\[
 \mu=\frac{25}{128},\quad \Lambda=\frac{153}{256},\quad
 v(s)=\min\left\{\frac{1+s}{2},\frac{57}{64}-\frac32s\right\},
 \qquad B(z)=\Lambda-\frac38z.
\]
For each nonempty $I\subseteq\{1,2,3,4\}$ define
\begin{align*}
 s_I(t)&=\sum_{i\in I}t_i,\\
 G_I(t)&=\max_{j\notin I}
       \left(s_I(t)+\sum_{\substack{k<j\\k\notin I}}t_k+2t_j\right),
       \qquad \max\varnothing=0.
\end{align*}
Include the endpoint $v(s_I(t))$ in the finite set $\mathcal C(t)$
if and only if
\begin{equation}\label{eq:four-prefix-subset-gate}
 v(s_I(t))>\max\{1/2,W(t),G_I(t)\}.
\end{equation}
Also include $B(z)$ if
\begin{equation}\label{eq:four-prefix-W-gate}
 x<\mu,\qquad B(z)>\max\{1/2,W(t),v(x)\}.
\end{equation}
Define $\mathcal E=\{t:\mathcal C(t)\ne\varnothing\}$ and
$L_4(t)=\max\mathcal C(t)$ on $\mathcal E$.  This is a finite
piecewise affine prescription, with strict eligibility gates.  All
its endpoints are at most $\Lambda<3/5$.

Each actual sieve level $\nu$ used for a selected endpoint must
satisfy
\[
 \max\{1/2,W(t),G_I(t)\}<\nu<v(s_I(t)),
\]
or the analogous interval between the two sides of
\eqref{eq:four-prefix-W-gate}.  An endpoint cap does not authorize
levels below its gate.  On a compact subregion of these strict
inequalities there is a fixed positive margin on every side.
Above $x=\mu$ only the explicitly tested selected-subproduct branches
are used.  In particular the full raw eight-branch envelope is not
extended beyond its proved largest-prime range.

\begin{lemma}[Averaged ordinary upper sieve with four prime slots]
\label{lem:four-prefix-upper-interface}
On any compact positive prefix box contained in a strict branch of
$\mathcal E$, choose $\nu$ between its gate and cap, with fixed
margins.  The ordinary residual upper-sieve weights of level
$N^{\nu-S+o(1)}$, sifted below the smallest specified prime, have
the averaged distribution remainder of
Lemma~\ref{lem:prefix-interface}, now with four specified primes.
Consequently their limiting upper coefficient at the prefix point
$t$ is
\begin{equation}\label{eq:four-prefix-ordinary-kernel}
 \mathcal O_\nu(t)
 =\frac4{\nu-S}\,
       \mathcal A\left(\frac{\nu-S}{w}\right),
 \qquad
 \mathcal A(s)=\frac{sF(s)}{2e^\gamma}\quad(s\ge2),
 \quad \mathcal A(s)=1\quad(0<s<2).
\end{equation}
The coefficient is in units of $\HN$ and is integrated against
$d\mu(t)$.  The assertion is averaged over the prefix box and does
not assert an estimate for each particular quadruple of primes.
\end{lemma}

\begin{proof}
The required change to Lemma~\ref{lem:prefix-interface} is at the
modulus atom, including its factorization; it is not a substitution
of $r=4$ in a theorem stated only for $r=1,2$.
Choose four disjoint ordered intervals
$I_i=(N^{\ell_i},N^{u_i}]$, a fixed level $\nu$, and
$T_+=\sum_i u_i$.  Make the intervals short enough that the upper
endpoint vectors satisfy every required square test with margin
$20\delta$, and that the entire chosen cap exceeds $\nu$ with
margin $20\delta$.  Set
\[
 R=N^{\nu-T_+-20\delta},\qquad
 Z=N^{u_*},\qquad
 u_* =\min\left\{\ell_4-\xi,
                 \frac{\nu-T_+-20\delta}{2+\rho}\right\}>0,
\]
where $\xi,\rho>0$ are fixed and sufficiently small.  Thus $Z$ is
below every specified prime and below the original quotient cutoff.
Lowering the cutoff gives a pointwise upper bound for the original
positive quotient count.  Also $\log R/\log Z\ge2+\rho$; no
linear-sieve formula with an inadmissible small parameter is invoked.

Construct the ordinary modified residual upper weights $\rho_R^+$
independently of the particular primes inside the four intervals.
Use the atoms of
\href{https://arxiv.org/html/2109.02851v2}{Lichtman 2021,
(5.14)--(5.18) and Theorem 2.10}, at residual level $R$.
Convolve each residual atom with
$\chi_{I_1}*\chi_{I_2}*\chi_{I_3}*\chi_{I_4}$.
Choose the construction precision $\eta$ small enough that its
presieving threshold $z_0=N^{c(\eta)}$ is below $Z$ and every
specified interval.  Every residual rough prime is below $Z$.
The specified slots are therefore exactly the largest four slots
in the concatenated rough vector.  There is no residual-prime
interleaving to be justified.

The four initial $\mathcal D^{\rm well}$ inequalities hold because
$W(t)<\nu$ with the prescribed margin.  For every later rough prime
$q_j$, residual well support gives
\[
 (p_1p_2p_3p_4)q_1\cdots q_{j-1}q_j^2
 \le N^{T_+}R^{1+O(\eta)}\le N^{\nu-19\delta}
\]
after decreasing $\eta$.  Thus the complete rough modulus is in
the required buffered $\mathcal D^{\rm well}$ support.  This check
also covers the residual atom with no rough slots: that atom still
has four specified rough slots, and is not reclassified as a
rank-zero modulus of level below $1/2$.

For a selected branch put $b_0=\prod_{i\in I}p_i$.
The inverse critical-interval identity in
Lemma~\ref{lem:selected-subproduct-factorization} applies because
$\nu<v(s_I)$ with a margin.  For each unselected specified slot its
remaining-square inequality is exactly the corresponding term in
$G_I<\nu$.  Before any subsequent residual prime all selected and
unselected specified slots have occurred, so its remaining-square
inequality is exactly the complete-vector inequality just checked.
The selected-subproduct lemma supplies a factorization for every
allowed Type II parameter $A$, including its $C<1$ case by
Lemma~\ref{lem:two-factor-greedy-repair}.  For the extra $B(z)$
branch use the repaired $W$ construction: on $x<\mu$ the inner
$W(x,y,z)$ cap is $B(z)$ and its gate is $v(x)$, as verified in
Lemma~\ref{lem:variable-envelope-simplification}.  Its hypotheses
and the complete well support have fixed margins as well.

Here are the coefficient details needed for this transfer.  The
factorization is performed on the upper endpoint vector, which
gives a fixed partition of its prime slots valid for every tuple
in the rough box.  The selected lemma explicitly supplies this
uniform partition on both sides of $C=1$.  Retain each original
scalar coefficient of the residual sieve atom, its signs, and its
reciprocal factorials for repeated residual interval labels.
No label crosses from a specified slot to a residual slot, because
$Z<\inf I_4$.  If a factor group contains $m_i$ rough slots,
divide its tuple convolution by $m_i!$.  The resulting coefficient
is $1$-bounded, even before removing nonsquarefree products; multiply
the original scalar by the product of these group factorials.
This is an exact convolution identity, and does not introduce a
second copy of the original repeated-label normalization.
The finite number of atoms and the sum of the absolute values of
their new scalars depend only on the fixed construction precision.

Combine the original signed small-prime weight $\beta$ with the
third factor.  Its support is $e\le N^{O(\eta)}$ on primes below
$z_0$, whereas that factor's rough primes exceed $z_0$.
Unique small/rough factorization proves that its coefficient stays
$1$-bounded.  Attaching $e$ to the third factor multiplies the three
factorization monomials by $e^2,e^2,e^{2+2\sigma}$, respectively;
the fixed slack pays these losses when $\eta\ll\delta$.
As in the existing prefix interface, allowing unrestricted prime
convolutions and then removing rough prime-square collisions costs
\[
 \ll (N+N^{\nu+O(\eta)})z_0^{-1}(\log N)^{C}
   =O_A(N(\log N)^{-A}).
\]
Thus neither squarefreeness nor signed atoms are replaced by an
unproved positivity assertion.

The distribution step is now precisely the existing atom argument.
The three factorization inequalities feed
\href{https://arxiv.org/html/2309.08522v1}{Lichtman 2023,
Proposition 5.2}, with $A_0=UX^{-\zeta}$ in the Type II range
on a local prime scale $X$.
The smallest saving after replacing $A_0$ by $U$ is
$(1-2\sigma)\zeta=9\zeta/16$.
The rough Siegel--Walfisz and bounded-complement steps, and the
one-long and two-long cases, are exactly those of the existing
prefix interface, with the corrected exact coprime mean of
Section~\ref{sec:one-long-coprime-mean}.
They only require a fixed number of normalized bounded atom factors
and total support below $X^{3/5}$, both just checked; they do not
require the number of specified slots to be at most two.

For clarity the moving residue is handled on the actual prime
scale.  Discard primes $P\le N^{1-\kappa}$, at cost
$N^{1-\kappa+o(1)}$ with the fixed divisor weights, and use dyadic
$P\asymp X$.  Then $\alpha=\log N/\log X=1+O(\kappa)$.
The strict cap and gate margins absorb this change and the
simultaneous rescaling of all exponent coordinates.  For the
selected branches this is the verified
$V_{\alpha,7/32}$ interval, not an unperturbed largest-prime
envelope.  The $W$ branch also retains $x<\mu$ with a fixed
margin before this change of scale.  Choose $\kappa$ small enough
that $N<X^{1+\epsilon}$; Proposition 5.2 explicitly permits this
moving residue.  This proves the uniform averaged remainder
\[
 \sum_{\substack{p_i\in I_i\\(p_1p_2p_3p_4,N)=1}}
 \sum_{\substack{q\mid P(Z)\\(q,N)=1}}
 \rho_R^+(q)E_N(p_1p_2p_3p_4q)
 =O_A(N(\log N)^{-A}).
\]

For the main term keep the full presieve factor $V_N(z_0)$, as in
the final part of Lemma~\ref{lem:prefix-interface}.  The one-sided
dimension-one estimate below $z_0$ is uniform in even $N$; above
$z_0$ the omitted primes dividing $N$ have total reciprocal weight
$o(1)$.  Consequently the main term is
\[
 \pi(N)V_N(Z)\bigl(F(\log R/\log Z)+o_\eta(1)\bigr)
 \sum_{\substack{p_i\in I_i\\(p_1p_2p_3p_4,N)=1}}
       \frac1{\varphi(p_1p_2p_3p_4)}.
\]
Here $V_N(Z)=(2e^{-\gamma}+o(1))C(N)/(u_*\log N)$ uniformly,
and the reciprocal prime sums tend to the product measure $d\mu$.
First let $N\to\infty$ at fixed margins and precision, then refine
the boxes and precision within those margins, and finally send
$\delta,\xi,\rho$ to zero.  The limiting cutoff exponent is
$\min\{w,(\nu-S)/2\}$.  If $(\nu-S)/w\ge2$, this gives
$2e^{-\gamma}F((\nu-S)/w)/w$; otherwise lowering gives exactly
$4/(\nu-S)$.  These are \eqref{eq:four-prefix-ordinary-kernel}.
\end{proof}

Define $\mathcal O(t)=\mathcal O_{L_4(t)}(t)$ on $\mathcal E$,
and $\mathcal O(t)=+\infty$ elsewhere.  The endpoint here is a limit
from strictly admissible levels.  Since $F$ is nonincreasing and
$4/(\nu-S)$ decreases in $\nu$, the largest eligible endpoint
gives the smallest ordinary option.  Moreover $W\ge S+w$, so
$L_4-S>w>a$ on $\mathcal E$; no singularity from zero residual
level occurs at its boundary.

\begin{proposition}[Four-prefix safe minimum]
\label{prop:four-prefix-safe-minimum}
Assume the existing local switching upper bound with coefficient
$K$ on positive prefix boxes of $\Omega_j$, for $j=10,11$.
Assume also $d+4h<1$, so the following Buchstab arguments exceed
one throughout the two domains.  Then a valid upper constant is
\begin{equation}\label{eq:four-prefix-safe-minimum}
 U_j^{\rm safe}
 =\int_{\Omega_j}
   \min\left\{\frac Kz\omega\left(\frac{1-S}{z}\right),
                   \mathcal O(t)\right\}\,d\mu(t),
 \qquad \Upsilon_j\le (U_j^{\rm safe}+o(1))\HN.
\end{equation}
Here $\Upsilon_{11}$ means its full fourth-prime range.  If
$a>1/12$, the ordinary option has the simpler exact form
\begin{equation}\label{eq:four-prefix-rational-kernel}
 \mathcal O(t)=\frac4{L_4(t)-S}\quad(t\in\mathcal E).
\end{equation}
This proposition supplies an analytic integral to be bounded; it
does not certify a numerical value of that integral.
\end{proposition}

\begin{proof}
Take a common finite positive partition for both upper bounds.
Inside a box separated from the strict cap, gate, and ordering
boundaries, select a branch and an admissible level with fixed
margins.  Both estimates bound the same averaged original
quotient count.  Select the smaller of the two box-level upper
bounds, then sum.  This does not take the minimum of unproved
estimates for individual prime quadruples.
Use switching alone on the remaining boxes, including those
meeting a boundary.  The finite gates and endpoint comparisons
are piecewise affine; their nontrivial equality sets have measure
zero.  Remove thin strips about these sets before refining the
partition.  On these bounded domains $t_i\ge a>0$ and
$(1-S)/z>1$, so the switching kernel and its measure are bounded.
The strips therefore contribute $o(1)$ as their widths tend to
zero, by the existing local switched estimate.  Refine within the
fixed margins, approach each eligible endpoint, and then remove
the margins.  This proves \eqref{eq:four-prefix-safe-minimum}.
The parameter inequality displayed in the proposition is a simple
sufficient condition for this kernel domain: $S+z<d+4h<1$.
Other constraints needed for an entire eleven-term Chen weight
remain separate and are not inferred from this local proposition.

To prove the simplification, let $k=|I|$ and use
$S\ge s_I+(4-k)a$ and $w\ge a$.  Whenever the selected endpoint
is eligible its residual exponent is positive, so
\[
 \frac{v(s_I)-S}{w}\le
 \begin{cases}
  1/(2a)-7/2,&k=1,\\
  1/(2a)-3,&k=2,\\
  57/(64a)-17/2,&k=3,\\
  57/(64a)-10,&k=4.
 \end{cases}
\]
The first two inequalities use $v(s)\le(1+s)/2$ and
$s_I\ge ka$; the last two use $v(s)\le57/64-3s/2$.
For the $B$ endpoint,
\[
 \frac{B(z)-S}{w}\le\frac\Lambda a-\frac{35}{8}.
\]
If $a>1/12$, all five upper bounds are strictly below $3$.
At $a=1/12$ their values are $5/2,3,35/16,11/16,179/64$,
respectively.  The normalized function $\mathcal A(s)$ is exactly
$1$ for $0<s\le3$, including the cutoff-lowering convention.
This proves \eqref{eq:four-prefix-rational-kernel}.  For smaller
$a$ the general $F$ kernel remains the justified formula; no
extrapolation of $\mathcal A=1$ is made.
\end{proof}

\paragraph{Why the domain is filtered.}
The four square tests cannot be dropped even below $x=\mu$.
For the original parameters $a=100/1149$, $h=50/309$,
$d=1/2-3a$, the interior point
$(x,y,z,w)=(.155,.150,.145,.140)\in\Omega_{10}$ has
$W=.730>\Lambda$.
The point $(.200,.140,.130,.120)\in\Omega_{11}$ has
$W=.710>\Lambda$.  Neither admits an endpoint of this construction.
They are handled by switching in \eqref{eq:four-prefix-safe-minimum}.
Conversely $(.170,.110,.100,.090)\in\Omega_{11}$ admits the
subset $I=\{3,4\}$: $s_I=.190$, $v(s_I)=.595$,
$G_I=.580$, and $W=.560$.  This demonstrates why all fifteen
subsets can be useful, even though only three specified slots
appeared in the earlier envelope.  Above the largest-prime cutoff,
$(.202,.092,.090,.088)\in\Omega_{11}$ admits $I=\{2,3\}$:
its cap is $.591$, its selected gate is $.586$, and $W=.560$.
All these terminating decimals are exact rational numbers.

\paragraph{Accounting for the earlier classical double-sieve option.}
Bounds from different methods may be retained as a third entry in
the same pointwise minimum when their common positive prefix class
has been proved.  Improvements from overlapping regions may not
be added.  In fact the previous small region
$100/1149<w<z<y<x<12/125$ is already improved by the ordinary
option here beyond its stated classical $H(2)$ option.  On it
$W<5(12/125)<1/2<v(x)$, $x<\mu$, and
$v(x)-1/2>a/2>H_0\Lambda$, where $H_0=223939/10^7$.
Since $v(x)-S\le\Lambda$ and $1/2-S>0$, this proves
\[
 \frac4{v(x)-S}
 <\frac{4(1-H_0)}{1/2-S}.
\]
Thus the previously certified gain exceeding $1/720000$ is
subsumed on its region and must not be added to a computation of
\eqref{eq:four-prefix-safe-minimum}.

The distribution and switching inputs used here are the established
derived interfaces and their cited published theorems.  The saved
JSON records the exact numerical domain and dependencies.  A
floating quadrature of this integral is a diagnostic only; a new
rounded Chen coefficient requires its own outward certificate.
